\centerline{\bf \Huge Вписанные углы, счёт дуг.}

\begin{flushright}
        \scriptsize{— Патрик? Это ты?\\
                    — Что?\\
                    — Я сказал «Это ты?»\\
                    — Я тебя не слышу. Здесь слишком темно.
}
\end{flushright}

\problem{1}
\textbf{Вспомогательные леммы.} На окружности выбраны точки $A, B, C$ и $D$ (именно в таком порядке по часовой стрелке)

\begin{enumerate}
    \item[(a)] Отрезки \( AC \) и \( BD \) пересекаются в точке \( P \). Докажите, что \( \angle APD = \frac{1}{2} (\overset{\frown}{AD} + \overset{\frown}{CB}) \)
    
    \item[(b)] Лучи \( AB \) и \( DC \) пересекаются в точке \( Q \). Докажите, что \( \angle AQD = \frac{1}{2} (\overset{\frown}{AD} - \overset{\frown}{CB}) \)
    
    \item[(c)] Докажите, что \( AD \parallel BC \) тогда и только тогда, когда \( \overset{\frown}{AB} = \overset{\frown}{CD} \)
\end{enumerate}

\problem{2}
Шестиугольник $ABCDEF$ вписан в окружность. Докажите, что
\[\angle A + \angle C + \angle E = \angle B + \angle D + \angle F\]

\problem{3}
Треугольник $ABC$ вписан в окружность $\omega$. Биссектрисы его внутренних углов B и C пересекают $\omega$ в точках $X$ и $Y$ . Докажите, что прямая $XY$ отсекает от угла $BAC$ равнобедренный треугольник.

\problem{4}
\textbf{Теорема Паскаля для самых маленьких.} Шестиугольник \( A_1 A_2 A_3 A_4 A_5 A_6 \) вписан в окружность. Оказалось, что \( A_1 A_2 \parallel A_4 A_5 \) и \( A_2 A_3 \parallel A_5 A_6 \). Докажите, что \( A_3 A_4 \parallel A_6 A_1 \).

\problem{5}
Четырехугольник $ABCD$ вписан в окружность. Докажите, что биссектрисы углов $ABD$ и $ACD$ пересекаются на этой окружности. Где именно они пересекутся?

\problem{6}
В окружность вписан четырёхугольник $ABCD$. Его вершины разбивают окружность на четыре дуги. Докажите, что отрезки, соединяющие середины противоположных дуг, перпендикулярны.

\problem{7}
Три равные окружности имеют общую точку $H$ и попарно повторно пересекаются в точках $A, B$ и $C$, причем точка H лежит внутри треугольника $ABC$. Докажите, что $H$ $-$ ортоцентр (то есть, точка пересечения высот) треугольника $ABC$.

\problem{8$^*$}
Окружность $\omega$ описана около остроугольного треугольника $ABC.$ Точка $H$ --- основание высоты $AH$ этого треугольника, а точка $G$ --- его точка пересечения. Луч $GH$ пересекает $\omega$ в точке $P.$ Докажите, что $\angle BPH = \angle ABH$
